% Source PDF page 79; manual page 2-50.
\subsection{Specific Load Factors}\label{sec:specific-load}

One definition of the specific-load vector is the vehicle's inertial acceleration minus
the acceleration caused by gravity:
\begin{equation}
  \vec a_{sp}=\vec a_I-\vec a_{grav}.
  \label{eq:specific-load-inertial-definition}
\end{equation}
An equivalent definition is the sum of the aerodynamic and propulsive forces divided
by vehicle mass:
\begin{equation}
  \vec a_{sp}
  =\frac{\sum\vec F_{aero}+\sum\vec F_{prop}}{m}.
  \label{eq:specific-load-force-definition}
\end{equation}
Specific-load factors are the components of $\vec a_{sp}$. They are used to establish
structural requirements and limits and are commonly referred to as vehicle
\emph{g loading}. The function \texttt{compute\_derivs} evaluates
\begin{longtable}{@{}>{\ttfamily}p{0.17\textwidth}p{0.21\textwidth}p{0.54\textwidth}@{}}
  nx & $\vec a_{sp}\mathbin{\cdot}\hat x_b$ &
    Specific-load component along the body $x$ axis, commonly called axial g's.\\
  ny & $\vec a_{sp}\mathbin{\cdot}\hat y_b$ &
    Specific-load component along the body $y$ axis.\\
  nz & $\vec a_{sp}\mathbin{\cdot}\hat z_b$ &
    Specific-load component along the body $z$ axis.\\
  ntotal & $\lVert\vec a_{sp}\rVert$ &
    Magnitude of the specific-load vector, commonly called total g's.\\
\end{longtable}
The normal, or lateral, specific load is
\begin{equation}
  a_{sp_n}
  =\sqrt{
    \left(\vec a_{sp}\mathbin{\cdot}\hat y_b\right)^2
    +\left(\vec a_{sp}\mathbin{\cdot}\hat z_b\right)^2
  }.
  \label{eq:normal-specific-load}
\end{equation}
